\section{Holonomic $\mathcal{D}$-modules}
Let $X = \text{spec} R$ be a smooth affine scheme with algebraic coordinates $(x_1, \ldots, x_n)$ and let $M$ be a finite generated $\mathcal{D}_X$-module. We already know $\text{ch}(M) \subset T^*X$ is a conic involutive (or coisotropic) subvariety. In particular, $\dim \text{ch}(M) \geq n$. 

\begin{definition}
    $M$ is called $\textbf{holonomic}$ if $\dim \text{ch}(M) = n$.
\end{definition}

\begin{corollary}
    If $M$ is a holonomic $\mathcal{D}_X$-module, then $\text{ch}(M) \subset T^*X$ is a Lagrangian subvariety, and $$CC(M) = \sum_{E_i \subset X\atop \text{irr subvarieties}} n_i T^*_{E_i} X$$
\end{corollary}

\begin{remark}
    Recall that Lagrangian means $W=W^{\perp}\iff \mathrm{dim}_{k}W=n$, and we define $V\subseteq T^{\ast}X$ is Lagrangian if for any smooth closed point $p\in V$, we have $\mathcal{T}_{V,p}\subseteq \mathcal{T}_{T^{\ast}X,p}$ is Lagrangian. Here $T^{\ast}_{E_{i}}X$ is the conormal bundle of irreducible subvariety $E_{i}$. 
\end{remark}

\begin{example}
    \quad
    \begin{enumerate}
        \item $R$ is a $\mathcal{D}_X$-module. $CC(R) = T^*_X X\simeq X \Rightarrow R$ is holonomic.
        \item Let $M$ be an integrable connection on $X$, we obtain that $ch(M)=T^{\ast}_{X}X$ and $\mathrm{dim}_{k}V(\xi_{1},\dots,\xi_{n})=n$, we have $M$ is holonomic. 
    \end{enumerate}
\end{example}

\section{The de Rham Complexes and Analytification}

Let $M$ be a $D_X$-module, and denote $\partial_i = \partial_{x_i}$ be the vector field along $x_i$.  
Define a complex $$K(M; \partial_1) = [M \xrightarrow{\partial_1} M].$$Here $[M\xrightarrow{\partial_{1}} M]$ means a complex concentrated on $\mathrm{deg}=0$ and $\mathrm{deg}=1$, with $\partial_{1}$ as coboundary map. 

\quad

\begin{note}[The cone of a morphism of complexes]\label{note 6.2}
Let $f: A^\bullet \longrightarrow B^\bullet$ be a morphism of complexes:

\begin{center}
 \begin{tikzcd}
 \cdots \arrow[r] & A^{n-1} \arrow[r, "d_A^{n-1}"] \arrow[d, "f^{n-1}"] & A^n \arrow[r, "d_A^n"] \arrow[d, "f^n"] & A^{n+1} \arrow[r] \arrow[d, "f^{n+1}"] & \cdots \\
 \cdots \arrow[r] & B^{n-1} \arrow[r, "d_B^{n-1}"] & B^n \arrow[r, "d_B^n"] & B^{n+1} \arrow[r] & \cdots
 \end{tikzcd}
\end{center}

then $\mathrm{Cone}(f)=C^{\bullet}$ defined by $C^n = A^n \oplus B^{n-1}$ with differential $$d_C^n = \begin{pmatrix}
 d_A^n & 0 \\
 f^n & -d_B^{n-1}
 \end{pmatrix}: A^n \oplus B^{n-1} \longrightarrow A^{n+1} \oplus B^n.$$The cone construction gives a short exact sequence of complexes:
\begin{center}
\begin{tikzcd}
0 \arrow[r] & B^\bullet[-1] \arrow[r] & Cone(f)^\bullet \arrow[r] & A^\bullet \arrow[r] & 0.
\end{tikzcd}
\end{center}
\end{note}


With Note \ref{note 6.2}, for $k \geq 2$, we inductively define     $$K(M; \partial_1, \cdots, \partial_k) = \text{cone}\left(K(M; \partial_1, \cdots, \partial_{k-1}) \xrightarrow{\partial_k} K(M; \partial_1, \cdots, \partial_{k-1})\right).$$ Then the $\textbf{de Rham complex}$ of $M$ is   $$K(M; \partial_1, \cdots, \partial_n),$$ denoted as $\mathrm{DR}(M).$

\begin{remark}
    In fact, it's Koszul complex with $(\partial_{1},\dots,\partial_{n})$ a finite sequence of central elements in $M$, defined in a mapping cone way.   
\end{remark}

Meanwhile, we have a coordinate-free way to define $\mathrm{DR}(M)$: $$M \to \Omega^{1}_{X} \otimes_{R} M \to \Omega^{2}_{X} \otimes_{R} M \to \cdots \to \Omega^{n}_{X} \otimes_{R} M.$$With $m\rightarrow \sum\limits_{i}dx_{i}\otimes \partial_{i}(m)$, and $\eta \otimes m \mapsto \sum dx_i \wedge \eta \otimes \partial_i(m)$ and so on...

\begin{remark}
    Here $\sum\limits_{i} dx_i \otimes \partial_i \in \Omega^1 \otimes_R \mathcal{T}_X$ is canonical i.e. independent of coordinates.
\end{remark}

\begin{lemma}
    We have $$\mathrm{DR}(M) \cong K(M; \partial_1, \ldots, \partial_n).$$The coordinate-free de Rham complex is isomorphic to the Koszul complex on the partial derivatives.
\end{lemma}
\begin{proof}
    When $n = 2$,  $\mathrm{DR}(M):0 \to M \to \Omega^{1}_{X} \otimes M \to \Omega^{2}_{X} \otimes M \to 0.$ $$K(M; \partial_{1}, \partial_{2})= M \to M \oplus M \to M\atop m \mapsto (\partial_{1} m, \partial_{2} m)\quad(m_1, m_2) \mapsto \partial_2(m_1) - \partial_1(m_2)$$On the other hand, we have $M\rightarrow \Omega^{1}_{X}\otimes_{R} M\rightarrow \Omega_{X}^{2}\otimes_{R}M$ to be $m\mapsto dx_{1}\otimes\partial_{1} m+dx_{2}\otimes\partial_{2}m$, $dx_{1}\otimes m_{1}+dx_{2}\otimes m_{2}\mapsto dx_{1}\wedge dx_{2}\otimes(\partial_{1} m_{2}-\partial_{2} m_{1})$
	
    The general case is similar. 
\end{proof}
	
\begin{example}
\quad
\begin{enumerate}
    \item Given $\alpha \in \mathbb{C}$, define $M_{\alpha} := \mathbb{C}[x, \frac{1}{x}] \cdot x^{\alpha}$, where $x^{\alpha}$ is a formal symbol such that $\partial_x x^{\alpha} = \alpha \frac{1}{x}\cdot x^{\alpha}$, then  $M_{\alpha}$ is a $\mathcal{D}_X$-module with: $$\partial_{x} \frac{1}{x^{k}} \cdot x^{\alpha} = (\alpha - k) \cdot \frac{1}{x^{k+1}} x^{\alpha} \quad \text{for } k \gg 0$$Since we obtain $\frac{1}{x^{k}}\cdot x^{\alpha}$ as a generator of $M_{\alpha}$ as $\mathcal{D}_{X}$-module, we have $M_{\alpha} \simeq \mathcal{D}_{\alpha} /\mathrm{Ann}_{\mathcal{D}_{\alpha}}( \frac{1}{x^{k}}\cdot x^{\alpha})\simeq\mathcal{D}_X / \mathcal{D}_X \cdot (x\partial_x - (\alpha - k))$. With $\mathrm{ch}(M)=V(\sigma(P))\subseteq T^{\ast}X$ in which $\sigma$ is the principal symbol, and $\sigma(P)=x\xi$, consider $T^{\ast}X\xrightarrow{\pi} X,\quad (x,\xi)\mapsto x$, we obtain $V(x)=\pi ^{-1}(0)=T^{\ast}_{\{0\}}X$ and $V(\xi)=T^{\ast}_{X}X$ zero section.  $\mathrm{ch}(M)=V(\sigma(x\partial_{x}))=T^{\ast}_{\{0\}}X\cup T_{X}^{\ast}X$ and  $$CC(M_{\alpha}) = T_{\{0\}}^*X + T_X^*X.$$Since $\mathrm{dim}ch(M)=1$ for $\mathrm{dim}T^{\ast}_{X}X=\mathrm{dim}X=1$, and $T^{\ast}_{\{0\}}X$ is the fiber of $T^{\ast}M$ at $0$, thus the cotangent space of $X$ at $0$, $\mathrm{dim}T^{\ast}_{\{0\}}X=\mathrm{dim}X=1$ as well. So  $M_{\alpha}$ is holonomic for all $\alpha \in \mathbb{C}$.
    \item Take $N = \mathbb{C}[\delta] \simeq \mathcal{D}_X / \mathcal{D}_X \cdot x \simeq \mathbb{C}[\partial_x]$, where $\delta$ is the Dirac delta function. Then $$CC(N) = T_{\{0\}}^*X$$also be holonomic. 

\end{enumerate}

\end{example}

Now, for $\alpha \notin \mathbb{Z}$, consider the de Rham complex of $M_{\alpha}$: $$M_{\alpha} \xrightarrow{\partial_1} M_{\alpha}$$

With $\text{Ker}(\partial_1) = \{f \cdot x^{\alpha} \mid \partial_1(f \cdot x^{\alpha}) = 0\}\simeq \{f \in \mathbb{C}[x, \frac{1}{x}] \mid xf' + \alpha f= 0\} = 0$, know that $x^{-\alpha}$ should be the solution of $xf'+\alpha f$, but $\alpha\not\in \mathbb{Z}$ makes $x^{-\alpha}\not\in \mathbb{C}\left[ x, \frac{1}{x} \right]$, thus has no solution in $\mathbb{C}[x, \frac{1}{x}]$. Yet, from ODE, we know that the differential equation $xf'+\alpha f=0$ has solutions at least locally, which suggests that the current definition of de Rham complex is not enough, and we need to extend scalars. 

\quad

Let $Y = \text{Spec } S$ be an affine variety of finite type over $\mathbb{C}$ with $\dim Y = n$. Define $Y \longrightarrow \mathbb{A}^{m}_{\mathbb{C}}$ by $$ \mathbb{C}[z_1, \ldots, z_m] \longrightarrow S\atop  z_i \mapsto \text{generators of } S$$And denote $I = I(Y) = \langle f_1, \ldots, f_k \rangle, \quad f_i \in \mathbb{C}[z_1, \ldots, z_m]$. 

In particular, $f_i$ are holomorphic functions on $\mathbb{C}^m$, then $f_1 = \cdots = f_k = 0$ defines an analytic subvariety $Y^{an}$ of the complex manifold $\mathbb{C}^m$. We have a continuous map: $$ \lambda: Y^{an} \longrightarrow Y $$For all$\ y \in Y^{\text{an}}$, by Implicit Function Theorem, $Y^{\text{an}}$ is smooth around $y$ if and only if the Jacobian matrix $J(f_1, \ldots, f_k)(y)$ at $y$ has rank $m-n$. 

\begin{definition} We define

\begin{enumerate}
    \item $Y$ is $\textbf{smooth}$ at a closed point $p$ if $$\dim_{\mathbb{C}} \mathfrak{m}_p/\mathfrak{m}_p^2 = n,$$ where $\mathfrak{m}_p$ is the maximal ideal of $p$.
    \item $Y$ is $\textbf{smooth}$ if $Y$ is smooth at every closed point.
\end{enumerate}
\end{definition}
\begin{lemma}
    $Y$ is smooth $\iff$ $Y^{\text{an}}$ is a complex manifold.
\end{lemma}
\begin{proof}
    Enough to prove for closed point $p$ of $Y$ we have: $$\text{rank}(J(f_1, \ldots, f_k)(p)) = m-n \quad (*)\iff \quad \dim_{\mathbb{C}} \mathfrak{m}_{p}/\mathfrak{m}_p^2 = n.$$Know that $p$ is also a closed point of $\mathbb{A}^{m}_{\mathbb{C}}$ with maximal ideal $\mathfrak{m} = (z_1 - \alpha_1, \cdots, z_m - \alpha_m)$. Define $$\mathbb{C}[z_1, \cdots, z_m] \to V = \mathbb{C}^m\atop f \quad \mapsto \quad \left( \frac{\partial f}{\partial z_1}(p), \cdots, \frac{\partial f}{\partial z_m}(p) \right)$$By HW, this induces an iso:$$
\mathfrak{\mathfrak{m}} /\mathfrak{m}^2 \xrightarrow[\varphi]{\simeq} V\simeq \Omega_{V,p}=\Omega_{V}\otimes R/\mathfrak{m}_{p}$$Since $I\subseteq \mathfrak{m}$ is generated by $f_{1},f_{2},\dots,f_{k}$, and for $h\in \mathbb{C}[z_{1},\dots,z_{m}]$, we can write $h=c+\sum(z_{i}-\alpha_{i})+\sum(z_{i}-\alpha_{i})(z_{j}-\alpha_{j})+\dots$, then we get $[hf_{i}]=c[f]$ in $\mathfrak{m} /\mathfrak{m}^{2}$ for $\mathfrak{m}^{2}$ makes only monomial with $\mathrm{deg}\leq 1$ matter and $f=\sum(z_{i}-\alpha_{i})+(\text{higher deg})$. So we have $\varphi(f_{1}),\dots,\varphi(f_{k})$ span $\varphi(I)=I+\mathfrak{m} /\mathfrak{m}^{2}$, and further, we obtain: $$
\dim(\varphi(I)) = \operatorname{rank} \left( J(f_1, \cdots, f_k)(p) \right).
$$
On the other hand, we have short exact sequence $$0\rightarrow\varphi(I)=I+\mathfrak{m}^{2} /\mathfrak{m}^{2}\rightarrow \mathfrak{m} /\mathfrak{m}^{2}\rightarrow \mathfrak{m}/\mathfrak{m}^{2}+I\rightarrow 0.$$
Know that $\mathfrak{m}_{p}=\mathfrak{m} /I$, $\mathfrak{m}_{p} /\mathfrak{m}_{p}^{2}=\frac{\mathfrak{m}/I}{\mathfrak{m}^{2}+I /I}\simeq \mathfrak{m} /\mathfrak{m}^{2}+I$,  with $\mathrm{dim}_{\mathbb{C}} \mathfrak{m} /\mathfrak{m}^{2}=m$, we get $\operatorname{rank} \left( J(f_1, \cdots, f_k)(p) \right)+\mathrm{dim}_{\mathbb{C}}\mathfrak{m}_{p} /\mathfrak{m}_{p}^{2}=m$, done. 
\end{proof}
\quad

Let $Z$ be a complex manifold (or variety) and $\mathcal{O}_{Z}$ the sheaf of all holomorphic functions on $Z$. 

We consider $Z=Y^{an}$, have $\lambda:Y^{an}\rightarrow Y$ induces a morphism of sheaves of rings $S \longrightarrow \mathcal{O}_{Y^{an}}$. If $M$ is an $S$-module, then $$
M^{\text{an}} \coloneqq M \otimes_S \mathcal{O}_{Y^{an}}
$$Then we finally define
$$
\operatorname{DR}(M) = K(M^{\text{an}}; \partial_1, \cdots, \partial_n)
$$
And $\mathrm{DR}(M)=[0 \rightarrow M^{an} \rightarrow M^{an} \otimes_R \Omega^1_X \rightarrow \cdots \rightarrow M^{an} \otimes_R \Omega^n_X \rightarrow0]=[0 \rightarrow M^{an} \rightarrow M^{an} \otimes_{\mathcal{O}} \Omega^1_{X^{an}} \rightarrow \cdots \rightarrow M^{an} \otimes_{\mathcal{O}} \Omega^n_{X^{an}}  \rightarrow0]$. 

Now, let's calculate $\mathrm{DR}(M_{\mathfrak{\alpha}})$. 

\begin{example}
    $X = \mathbb{A}^1_{\mathbb{C}}$,  $M_{\alpha} := \mathbb{C}[x, \frac{1}{x}] \cdot x^{\alpha}$.
    \begin{enumerate} 
    \item $\alpha \not\in \mathbb{Z}$, $\mathrm{DR}(M_{\alpha}) \cong [M_{\alpha} \xrightarrow{d_{\alpha}} M_{\alpha}]$. 
$\text{Ker}\partial_{x} \mid_{x \backslash \{0\}} \\ = \mathbb{C} \cdot x^{-\alpha} \cdot x^{d}$ since $\partial_{x}(x^{-\alpha}\cdot x^{\alpha})=\partial_{x}\cdot 1=0$. This implies $\text{Ker}\partial_{x} \mid_{x \backslash \{0\}}$ is the (non-trivial) local system $L_{\lambda}$ given by the multivalued function $x^{-\alpha}$, i.e. $T \cdot x^{-\alpha} = T \cdot e^{-\alpha \log x}= e^{-\alpha(\log x + 2\pi i)} = e^{-\alpha \cdot 2\pi i} e^{-\alpha \log x}= \lambda e^{-\alpha \log x}$, here $\lambda=e^{-\alpha2\pi i}$. 
$\partial_{x} \mid_{x \backslash \{0\}}$ is surjective, and then $\mathrm{DR}(M_{\alpha}) \mid_{x \backslash \{0\}} \simeq L_{\lambda}$, and  $H^{0}\mathrm{DR}(M_{\alpha}) \mid_{\{0\}} = H^{1}DR(M_{\alpha}) \mid_{\{0\}} = 0$. 
\item $\alpha \in \mathbb{Z} \Rightarrow \lambda = 1$, $\mathrm{DR}(M_{\alpha}) \mid_{x \backslash \{0\}} = \mathbb{C}_{x \backslash \{0\}}$, $H^{0}\mathrm{DR}(M_{\alpha}) \cong H^{1}\mathrm{DR}(M_{\alpha}) \cong \mathbb{C}_{\{0\}}$
\item $N=\mathbb{C}[\delta]\simeq \mathbb{C}[\partial_{x}]$, $\mathrm{DR}(N)\simeq \mathbb{C}_{\{0\}}[-1]$. 
    \end{enumerate}
\end{example}

\section{Basic version Riemman-Hilbert correspondence}

Let $Y$ be a complex manifold, $\mathcal{O}_Y$: the sheaf of holomorphic functions on $Y$, $\mathcal{L}$: a sheaf of $\mathbb{C}$-vector spaces. 

\begin{definition}
    $\mathcal{L}$ is called a $\textbf{local system}$ if $\forall y \in Y$, $\exists$ open neighborhood $U_y$ of $y$ s.t. $\mathcal{L}|_{U_y} \simeq\mathbb{C}^r_{U_y}$, where $\mathbb{C}^{r}_{U_y}$ is the constant sheaf on $U_y$.
\end{definition}

\begin{theorem}
    Let $X = \text{Spec } R$ be a smooth affine algebraic variety over $\mathbb{C}$ and $M$ an integrable connection on $X$, then $$M^{\text{an}} \simeq \mathcal{L} \otimes \mathcal{O}_{X^{\text{an}}},$$with $\mathcal{L}$ a local system.
\end{theorem}
\begin{proof}
   \quad

    Goal: $M$ is locally free, we want to show for small enough neighborhood, we can pick a basis $s_{1},\dots,s_{r}$ makes $\partial_{x_{i}}s_{j}=0$, thus
	We have known $M$ is a locally free $R$-module, $\forall x \in X$, take an open neighborhood $U \subset X$ s.t. $M|_U$ is free, then $M^{\text{an}}|_{U^{\text{an}}}$ is a free $\mathcal{O}_{U^{\text{an}}}$-module, generated by $\underline{\mathfrak{m}} = (m_1, \cdots, m_r)$.

	Denote $(x_1, \cdots, x_n)$ coordinates of $U^{\text{an}}$, then we have $\partial_{n} \underline{\mathfrak{m}} = X \cdot\underline{ \mathfrak{m}}$, where $X$ is a matrix of holomorphic functions on $U^{an}$. Expand $X$ w.r.t. $x_n$ in a small enough neighborhood of $x\in V\subseteq U^{an}$, we have $$X = \sum_{i\ge0} x_n^i \cdot X_i,$$where $X_i$ is a matrix of holomorphic functions in $(x_1, \cdots, x_{n-1})$. 

	Take $\Psi=\sum\limits_{i\ge 0} \frac{1}{i+1} x_{n}^{i+1}X_{i}$, so $\partial_{x}(\Psi)=X$, then $\Phi=\mathrm{\exp}(-\Psi)$, which is an invertible matrix. Set $\underline{S}=(s_{1},\dots,s_{r})=\Phi(\underline{\mathfrak{m}})$, so $\partial_{x}\underline{S}=\partial_{n}(\Phi)\cdot \underline{\mathfrak{m}}+\Phi \cdot(\partial_{n} \underline{\mathfrak{m}})=(\partial_{n}(\Phi)+\Phi \cdot X)\underline{\mathfrak{m}}=0$. 

	Since $\Phi$ is invertible matrix, we obtain $s_{1},\dots,s_{r}$ also a basis of $U^{an}|_{V}$, so we have $M^{an}|_{V}=\bigoplus\limits_{i=1}^{r}\mathcal{O}_{V}s_{i}$ s.t. $\partial_{n}s_{i}=0$. 

	Want to clarify the relation of $s_{i}$ and other $\partial_{x_{j}},j<n$, then Take $M_1 = \ker[M^{\text{an}}|_V \xrightarrow{\partial_n} M^{\text{an}}|_V]=\left\{ \partial_{n}\left( \sum f_{i}s_{i} \right)=0:f_{i}\in \mathcal{O}_{V} \right\}=\left\{ \sum\limits_{i=1}^{n}\partial_{n}(f_{i})s_{i}=0 \right\}.$ So we get $\sum_{i=1}^r \partial_n(f_i) s_i = 0$ for all $i=1,\dots,r$, i.e. $f_{i}$ is independent of $x_{n}$. 

	So we have $M_{1}=\bigoplus\limits_{i=1}^{r}\mathcal{O}_{V}^{n-1}s_{i}$, $\mathcal{O}_{V}^{n-1}$ are holomorphic functions that are independent of $x_{n}$, which suggests that we do not need the var $x_{n}$. And do the same thing for $M_{1}$,we get a new basis with $\partial_{n}s_{i}=\partial_{n-1}s_{i}=0$ for all $i=1,\dots,r$, and the change of basis is via an invertible matrix of homomorphic functions in $(x_{1},\dots,x_{n-1})$. 

	Continue the process, by induction, we finally obtain:  $$ M_n = \bigoplus_{i=1}^r \mathbb{C} \cdot s_i $$with $\partial_j(s_i) = 0 \quad \forall i, j$. And since every change of basis is by an invertible matrix, we know that for a small enough neighborhood $V \subseteq U^{an}$, we get $M^{an}|_{V}=\bigoplus\limits_{i=1}^{r}\mathcal{O}_{V}s_{i}$, with $\partial_{j}(s_{i})=0$ . 
\end{proof}

\begin{corollary}
    $\mathrm{DR}(M)|_{V}\stackrel{\text{quasi-isomorphism}}{\simeq}H^{0}\mathrm{DR}(M)|_{V}\simeq \mathbb{C}^{r}_{V}$. 
\end{corollary}
\begin{example}[Algebraic vs Analytic]
    Let $M = \mathbb{C}[x, \frac{1}{x}]$, $M' = \mathbb{C}[x, \frac{1}{x}] \cdot e^{\frac{1}{x}}$. 
    
    Know that $M$ and $M'$ are different as algebraic $D$-modules, but $M^{\text{an}}|_{\mathbb{C}^*} = \mathcal{O}_{\mathbb{C}^*}$, $M'^{an}|_{\mathbb{C}^{\ast}}=\mathcal{O}_{\mathbb{C}^{\ast}}\cdot e^{1/x}=\mathcal{O}_{\mathbb{C}^{\ast}}$, where $\mathbb{C}^{\ast}$ is the punctured complex plane. So $M^{\text{an}}|_{\mathbb{C}^*} \cong M'^{\text{an}}|_{\mathbb{C}^*}$ as analytic $\mathcal{D}$-modules and $DR(M^{\text{an}}) = DR(M'^{\text{an}}) \cong \mathbb{C}_{\mathbb{C}^*}^*$ Moreover, $M^{\text{an}} = \mathcal{O}_{\mathbb{C}}\cdot \frac{1}{x},M'^{an}=\mathcal{O}_{\mathbb{C}}\cdot e^{1/x}$ so $M^{an}$ and $M'^{an}$ are different as analytic $\mathcal{D}$-module. 
\end{example}

\section{Regular holonomic $\mathcal{D}$-module}

Let $X = \text{Spec } R$ be a smooth affine algebraic variety over $\mathbb{C}$ and $M$ a holonomic $\mathcal{D}_X$-mod. 
For any point $x \in X^{\text{an}}$, we consider the embedding:$$
i: \Delta \hookrightarrow X^{\text{an}}, \quad 0 \mapsto x.
$$where $\Delta$ is a small disk in the complex plane.
Then we have a morphism of sheaves of rings:$$R \to \mathcal{O}_{\Delta}, \quad f \mapsto f|_{\Delta}.$$We define $i^*M = M \otimes_R \mathcal{O}_{\Delta}$. 
\begin{definition}
    $M$ is called $\textbf{regular holonomic}$, if $i^{\ast}M$ is generated by $M_{\alpha}^{an}$ and $N^{an}|_{\Delta}$ in $\mathbf{K}_{0}(\mathcal{D}_{\Delta})$ for every small enough disk embedding $i:\Delta \hookrightarrow X^{an}$. 
\end{definition}

\begin{example}
    $M=\mathbb{C}\left[ x, \frac{1}{x} \right]\cdot e^{1/x}$ is not regular because $M$ is simple and $M\simeq \mathcal{D}_{X} /\mathcal{D}_{X}(x^{2}\partial_{x}+1)$. 
\end{example}